\head{第九版新核查：第0章20.1--20.5}
\source{EGA IV第一部分，分册印刷第116--136页，上传PDF第115--135页；止于20.5.15。作者勘误据IV第四部分，印刷第345、347--348页，上传PDF第344、346--347页}
本轮从§20章首说明接续，完整译入20.1“导子与代数扩张”、20.2“导子的函子性质”、20.3“拓扑环上的连续导子”、20.4“主部与微分”和20.5“微分模的基本函子性质”。共49个原著编号条目，保留57个原编号公式或交换图标签。原文没有编号的中间式没有擅加原著编号；为适应默认版心，少数长交换图分作相接的两段。

新增正文采用核实后的陈述和证明，原文错误及依据放在脚注。正式勘误、译者的数学修补、单纯记号或引用调整，分别记录。本轮没有重新逐句审定此前的全部译文与历史报告。下一接续点为20.6章首及（20.6.1）；已读取的20.6开头不计入已译正文。

\section*{本轮结论概览}
\addcontentsline{toc}{section}{本轮结论概览}
\begin{longtable}{@{}p{.17\linewidth}p{.19\linewidth}p{.55\linewidth}@{}}
\toprule 标识 & 原条目 & 判定与正文处理\\\midrule\endhead
I9-E01 & 20.4.13(ii) & 原微分模同构按字面不成立。改为沿增广张量到基环后的同构；给出特征2反例与一般表示。\\
I9-E02 & 20.5.4.2 & 一般交换方块中的导子图源对象错误。改用 $B\to B'$ 的自然性图，不改变微分模的典范映射。\\
I9-E03 & 20.3.2 & 若只要求满同态连续，导子下降后未必连续。正文明确商拓扑充分条件，并给出缺此条件的检验例。\\
I9-C01--C05 & 20.1.2；20.2.1；20.4.6--20.4.13 & 作者已有勘误：标量环、微分符号、相应证明等式、去掉整闭性。另收录20.5.4的上标修正。\\
I9-T01--T08 & 各相应条目 & 加法性、开理想、复合次序、源靶及交叉引用的明确化或修正，不统称为八个错误定理。\\
I9-N01--N05 & 各相应条目 & 双模作用、群结构、连接同态加法性、分裂张量证明及交叉幂等元的说明。属于补明步骤。\\
\bottomrule
\end{longtable}

\audit{I9-E01}{平方零扩张的微分模：原同构缺少沿增广取商}{原断言不成立；有具体反例及修订后的完整证明}
\oldauditsrc{（$0_{\mathrm{IV}}$，20.4.13，ii），分册印刷第127页／上传PDF第126页}
原文设 $L$ 为 $A$-模，$B=\mathrm D_A(L)=A\oplus L$，声称典范映射 $L\to\Omega^1_{B/A}$，$x\mapsto dx$，是双射；随后又把所用导子称为“$B$-导子”。这里有两个不同问题。任何 $B$-导子 $B\to M$ 都为零；即使把它改为 $A$-导子，原来的微分模同构仍然不成立。

\lead{一个二维计算足以反驳原式。}
令 $k=\mathbb F_2$，$L=k\epsilon$，$B=k[\epsilon]/(\epsilon^2)$。由多项式代数的微分模与商代数关系，
\[
 \Omega^1_{B/k}
 =B\,d\epsilon\big/\bigl(d(\epsilon^2)\bigr)
 =B\,d\epsilon\big/(2\epsilon\,d\epsilon)
 =B\,d\epsilon.
\]
因此 $\Omega^1_{B/k}$ 作为 $k$-向量空间有基 $d\epsilon,\epsilon d\epsilon$，维数为2；而 $L$ 的维数为1。映射 $a\epsilon\mapsto a\,d\epsilon$ 不满射。

也可以不预先使用商代数微分公式。对任意 $m\in B$，令
\[
 D_m(a+b\epsilon)=b\,m\quad(a,b\in k).
\]
在特征2下，
\[
 (a+b\epsilon)D_m(c+d\epsilon)
 +(c+d\epsilon)D_m(a+b\epsilon)=(ad+bc)m,
\]
恰等于 $D_m((a+b\epsilon)(c+d\epsilon))$。故每个 $m\in B$ 都能作为一个导子的 $D_m(\epsilon)$；这说明代表导子函子的微分模必须容纳自由的 $B$-生成元 $d\epsilon$。特别取 $m=1$，相应的 $B$-线性映射 $\Omega^1_{B/k}\to B$ 把 $\epsilon d\epsilon$ 映到 $\epsilon\ne0$，从而直接证明 $\epsilon d\epsilon\ne0$。

\lead{正文采用的正确陈述。}
设 $\varepsilon:B=A\oplus L\to A$ 为增广，则
\[
 L\xrightarrow{\sim}\Omega^1_{B/A}\otimes_{B,\varepsilon}A,
 \qquad x\longmapsto dx\otimes1.
\]
这不是额外要求 $2$ 可逆，也不要求 $L$ 有限生成；它对任意交换环 $A$ 及任意 $A$-模 $L$ 成立。

\begin{proof}
对任意 $A$-模 $M$，通过增广把它视为 $B$-模。于是 $L$ 在 $M$ 上作用为零。每个 $A$-导子 $D:B\to M$ 在 $A$ 上为零，由其在 $L$ 上的限制 $h$ 唯一决定；$h$ 是 $A$-线性的。反之，任意 $h\in\Hom_A(L,M)$ 定义
\[
 D(a,x)=h(x).
\]
由于 $(a,x)(b,y)=(ab,ay+bx)$，有
\[
 D((a,x)(b,y))=a h(y)+b h(x)
 =(a,x)D(b,y)+(b,y)D(a,x).
\]
所以这是导子，得到函子同构
\[
 \Hom_A(L,M)\simeq\Der_A(B,M)
 \simeq\Hom_A(\Omega^1_{B/A}\otimes_B A,M).
\]
此同构由 $x\mapsto dx\otimes1$ 诱导，故该映射双射。
\end{proof}

\lead{完整微分模的表示。}
更精确地，令
\[
 Q=(B\otimes_A L)/R,
\]
其中 $R$ 为由 $x\otimes y+y\otimes x$（$x,y\in L\subset B$）生成的 $B$-子模。则
\[
 \Omega^1_{B/A}\simeq Q.
\]
证明如下。给定 $B$-模 $M$，$A$-导子 $D:B\to M$ 相当于 $A$-线性映射 $h:L\to M$，另满足
\[
 xh(y)+yh(x)=0\quad(x,y\in L),
\]
因为 $xy=0$；这正是从 $B\otimes_A L$ 到 $M$ 的 $B$-线性映射零化 $R$ 的条件。因此 $Q$ 具有微分模的泛性质。再沿增广与 $A$ 张量，所有这些关系的系数 $x,y$ 均变成零，得到前述同构。这个表示同时解释了原文为何在一般 $B$-模目标上漏掉了条件。

\lead{验证脚本的作用。}
工程附带的有限域脚本枚举 $\mathbb F_2[\epsilon]/(\epsilon^2)$ 的4个元素，逐项核查每个 $D_m$ 的加法性、Leibniz公式及 $D_1(\epsilon)=1$。程序只是这个有限例子的可复算核对；任意 $A,L$ 的结论由上面的代数证明承担，不以有限枚举代替一般证明。

\audit{I9-E02}{一般基环方块中的微分自然性图}{源对象与线性类别错误；典范微分模映射本身仍成立}
\oldauditsrc{（20.5.4.2），分册印刷第130页／上传PDF第129页}
给定交换环同态图
\[
\begin{tikzcd}
 B\arrow[r,"v"]&B'\\
 A\arrow[r,"u"']\arrow[u]&A'\arrow[u]
\end{tikzcd}
\]
总有典范 $B'$-线性映射
\[
 \theta:\Omega^1_{B/A}\otimes_B B'\longrightarrow\Omega^1_{B'/A'}.
\]
原文的（20.5.4.2）下行却为 $B'\xrightarrow{1_{B'}}B'$，左竖箭头记作 $d_{B/A}\otimes1$，并称为 $A'$-同态图。在一般方块下，$d_{B/A}$ 的定义域是 $B$；不能不经构造就把它张量成一个定义在整个 $B'$ 上的导子。

\lead{直接检验。}
取 $A=B=A'=k$、$B'=k[t]$，$u=1_k$，其余为结构嵌入。此时
\[
 \Omega^1_{B/A}\otimes_B B'=0,
 \qquad\Omega^1_{B'/A'}=k[t]\,dt.
\]
原图的左竖箭头只能为零；若右下角仍为 $B'$、下行仍为恒等，则交换性将给出 $d_{k[t]/k}=0$。但求导 $f(t)\mapsto f'(t)$ 是 $k$-导子，且将 $t$ 映为1，故 $dt\ne0$。这反驳原图，而不是反驳 $\theta$ 的存在性。

\lead{正文的正确方块。}
\[
\begin{tikzcd}[column sep=huge]
 \Omega^1_{B/A}\otimes_B B'\arrow[r,"\theta"]&\Omega^1_{B'/A'}\\
 B\arrow[u,"{b\mapsto db\otimes1}"]\arrow[r,"v"']&B'\arrow[u,"d_{B'/A'}"'].
\end{tikzcd}
\]
这是 $A$-线性映射的交换图，交换性逐元表达为
\[
 \theta(db\otimes1)=d_{B'/A'}(v(b)).
\]
它正是（20.5.2）、（20.5.3）构造的复合。除非另给出结构，不能把左下角 $B$ 称为 $A'$-模。

若额外假设 $B'=B\otimes_A A'$，则确可定义
\[
 D'(b\otimes a')=db\otimes a',
\]
并核对它是 $A'$-导子。这个特殊情形是下一命题的场景，不能提前当成（20.5.4）的一般假设。作者正式勘误还补了同页目标微分模的上标 $1$；那项是排印修正，不包含本报告指出的源对象修正。

\audit{I9-E03}{连续导子的下降需要拓扑条件}{原文只要求连续时不足；正文明确商拓扑充分版本}
\oldauditsrc{（20.3.2）对（20.1.4）的连续推广，分册印刷第122页／上传PDF第121页}
代数版本给出扩张 $0\to L\to E\xrightarrow p C\to0$ 的自等价与 $\Der_A(C,L)$ 的对应。连续版本中，从连续映射 $D':E\to L$、$D'|_L=0$ 推出下降后的 $D:C\to L$ 连续，必须核查 $C$ 的拓扑。原文在本处仅说明“当然要假设 $p$ 连续”。本轮对其按字面的范围作检验，不把未写出的商拓扑假设静默加入原陈述。

\lead{反例的全部对象。}
取域 $k$，令
\[
 C=k[x_1,x_2,\ldots],\qquad
 I_n=(x_i\mid i\ge n).
\]
以 $(I_n)_{n\ge1}$ 为 $C$ 的零邻域基。这是分离的线性环拓扑：各 $I_n$ 是双边理想，且任意非零多项式只涉及有限多个变量，故不属于所有 $I_n$。

令 $\varepsilon:C\to k$ 为把所有变量映到零的增广，并令 $L=k$ 为通过 $\varepsilon$ 定义的离散 $C$-双模。作用连续，因为 $I_1=\Ker\varepsilon$ 是开理想且零化 $L$。取平方零扩张
\[
 E=\mathrm D_C(L)=C\oplus k,
 \quad(c,l)(c',l')=(cc',\varepsilon(c)l'+\varepsilon(c')l),
\]
但给 $E$ 赋予离散拓扑。$p:E\to C$ 是连续满同态，$L$ 的子空间拓扑也是离散的；然而 $C$ 的上述尾变量拓扑不是 $p$ 的商拓扑，后者应为离散拓扑。

\lead{不连续导子对应连续自等价。}
定义 $D:C\to k$ 为多项式所有一次项系数之和。由于每个多项式只有有限项，此定义有意义；$D(x_i)=1$，且
\[
 D(fg)=\varepsilon(f)D(g)+\varepsilon(g)D(f).
\]
这可由乘积的一次项逐项计算得到，故 $D$ 是 $k$-导子。但每个 $I_n$ 都含 $x_n$，且 $D(x_n)=1$，所以 $\Ker D$ 不包含任何 $I_n$；目标离散，因而 $D$ 不连续。

另一方面，
\[
 F:E\longrightarrow E,\qquad F(c,l)=(c,l+D(c))
\]
是固定商 $C$ 和核 $L$ 的 $k$-扩张自等价，逆为 $(c,l)\mapsto(c,l-D(c))$。两个映射在离散空间 $E$ 上都连续。因此只假设 $p$ 连续时，连续自等价不一定来自连续导子。代数对应的唯一性还保证不存在另一个连续导子产生同一个 $F$。

\lead{正文采用的充分条件与证明。}
正文要求：$L$ 带 $E$ 的子空间拓扑，$C$ 带 $p$ 诱导的商拓扑。设 $F$ 为连续自等价，则 $D'=F-1_E$ 连续，值落在 $L$，故作为到子空间 $L$ 的映射仍连续。它在 $L$ 上为零，下降为唯一 $D:C\to L$。由商拓扑的定义，$D\circ p=D'$ 连续蕴涵 $D$ 连续。反之，连续导子 $D$ 给出连续 $F=1_E+D p$，其逆 $1_E-Dp$ 也连续。于是连续版本的群同构成立。

本项只说明原文字面所列条件不足，以及上述明确版本充分；不宣称商拓扑对每个个别扩张都是必要条件。若采用的范畴从定义上只容许严格的拓扑扩张，则应把这一约定写明，反例即不在该较小范畴中。§18.5的 $\Exantop$ 则是离散商扩张群的归纳极限，原书还提醒它不直接等同于一个拓扑扩张集合；不能仅凭该记号替本处补上隐含拓扑约定。

\audit{I9-C01--C05}{作者正式勘误的落实}{五项文字修订及一项排印修订；不是六个新错误定理}
\oldauditsrc{勘误清单三Err$_{\mathrm{IV}}$10--14，分册第347--348页；A类20.5.4，第345页}
\begin{enumerate}[label=\arabic*)]
\item （20.1.2）末：将 $a\in A$、$aD$、“$A$-模”换为 $b\in B$、$bD$、“$B$-模”。这是自然标量结构的加强；原来限制标量到 $A$ 的说法本身不假。
\item （20.2.1）：系数模同态诱导的导子映射为 $B$-线性；限制基环得到的映射也为 $B$-线性，而不是仅用关于基环的双同态来描述。
\item （20.4.6.1）：外微分统一为 $d=p_2-p_1$；微分 $dx$ 为 $1\otimes x-x\otimes1$ 的类。
\item （20.4.8）证明：与上一约定一致的关系为 $p_1=p_2-d$。不能只改定义，不改后续等式。
\item （20.4.13，iv）：删除 $A$ 整闭的额外条件。若 $A\subseteq B$ 为整环，$B$ 整于 $A$，分式域扩张可分，则 $\Omega^1_{B/A}$ 仍是挠 $B$-模。用最小多项式乘上 $0\ne a\in A$ 清分母，得 $af'(x)dx=0$，且 $af'(x)\ne0$。
\item （20.5.4）：目标微分模记号补上遗漏的上标 $1$。它与I9-E02中的一般自然性图问题分开记录。
\end{enumerate}
所有项目直接采用作者修订后的正文，脚注说明初版差异。书末增补这六条勘误的选译；并未声称勘误清单三其余部分已经全部译完。

\audit{I9-T01--T08}{附加条件、复合次序和引文校订}{分别定位；不把所有修订计作命题错误}
\lead{T01：加法性。}
（20.1.1）末句的逆向表述，明确说“一个加法映射 $D$”。导子定义本来要求加法性；原证明中的 $D=v-v_0$ 已有此性质。正文的明确化防止把Leibniz等式单独当作导子定义。

\lead{T02：扩张同态的引文号。}
（20.2.2）实际使用的三种映射是改变系数模、拉回扩张和限制基环；其定义分别在18.3.6.1、18.3.6.2、18.3.6.4。原文指向18.3.1、18.3.3、18.3.4中的编号不能定位这些映射。正文按源靶改引，并没有更换函子本身。

\lead{T03：结构同态少了复合。}
在 $A\xrightarrow u B\xrightarrow v C$ 的情形中，$\mathrm D_C(L)$ 的 $A$-结构为
\[
 t\longmapsto(v(u(t)),0).
\]
原文第120页写成 $(u(t),0)$，第一分量还在 $B$ 中。正文补上 $v$；数学意义为恢复原结构同态。

\lead{T04：离散商必须来自开理想。}
（20.3.5）指标集合中的 $\mathfrak j\subset A$、$\mathfrak k\subset B$ 明确要求为开双边理想；这与原证明使用离散商环一致。对连续导子取核中的开理想后，还要与零化 $L$ 的开理想相交，才能把 $L$ 看成商环双模。此缩小不损失连续性，也不改极限范围。

\lead{T05：20.3.8的假设定位。}
其两条序列分别应在20.3.6和20.3.7的假设下讨论；原引20.3.5、20.3.6，前一个没有所需的环 $C$，后一个未采用交换代数限制。正文按实际序列修正。

\lead{T06：形式平展推论中的箭头。}
（20.5.8）的源为 $\Omega^1_{B/A}\otimes_B C$，故引用20.5.2.4；原引20.5.3.3的源为 $\Omega^1_{C/A}$，是另一条箭头。形式平展性使第一条箭头双射，结论不变。

\lead{T07：$S^{-1}A$ 的定义域。}
（20.5.9.1）允许任意乘法集 $S\subseteq B$，结论保持。后面的公式20.5.9.2还写 $S^{-1}A$，故正文明确这个附加公式中的 $S$ 取自 $A$，其在 $B$ 中的像继续记为 $S$。不能拿 $B$ 中任意元素直接作为 $A$ 的分母。

\lead{T08：20.5.12.3后的复合。}
原定义18.3.8.1给出
\[
 \eta:\Hom_C(\mathfrak k/\mathfrak k^2,L)
 \xrightarrow{\sim}\Exalcom_B(C,L),
\]
而连接映射为 $\partial:\Der_A(B,L)\to\Exalcom_B(C,L)$。因此从导子到余法向模的 $\Hom$ 的映射应为
\[
 \varphi=\eta^{-1}\circ\partial,
\]
而不是原文的 $\partial\circ\eta^{-1}$。后一复合源靶不合。若将 $D$ 对应到 $\partial(D)$，按构造，$\eta^{-1}(\partial(D))$ 就是 $x\bmod\mathfrak k^2\mapsto D(x)$；这也逐元验证它是预复合 $\delta_{C/B/A}$ 的同态。另一映射 $\psi=v^1\circ\eta$ 保持原方向。

\audit{I9-N01--N05}{证明步骤和结构数据的补明}{补证与澄清，不作为原断言反例}
\lead{N01：20.1.1不假设 $p$ 满射。}
该条只给出 $p:E\to C$ 的平方零核和一个提升 $v_0:B\to E$。所需的是核上的 $B$-双模结构：
\[
 b\cdot z=v_0(b)z,\qquad z\cdot b=zv_0(b).
\]
若换另一提升，其与 $v_0$ 之差属于核，乘以 $z$ 后为零。因此作用与提升无关。不必也不能在未假设 $p$ 满射时，将核当然地赋予整个 $C$ 的作用。正文与脚注明确这一点，避免校订者自己引入多余假设。

\lead{N02：扩张自等价的群结构。}
以 $F_D=1_E+D p$ 记相应自等价。因 $pD=0$，有
\[
 F_{D_1}F_{D_2}=1_E+(D_1+D_2)p,
 \qquad F_D^{-1}=1_E-Dp.
\]
这核实20.1.4中的群同构，而非仅集合双射。

\lead{N03：连接映射的加法性。}
将 $D_1,D_2$ 对应的扩张先沿商 $C$ 取纤维积，核变为 $L\oplus L$，结构映射的核分量为 $(D_1,D_2)$。再沿加法 $L\oplus L\to L$ 推出，核分量变为 $D_1+D_2$。故 $\partial(D_1+D_2)=\partial(D_1)+\partial(D_2)$。这补明了原文五处正合性之外，连接映射确为群同态的步骤。

\lead{N04：任意基变换下的核。}
（20.5.5）的结论不需要 $A'$ 对 $A$ 平坦。关键不是 $p:B\otimes_A B\to B$ 满射，而是它有 $A$-线性截面 $j_1$。因此
\[
 0\longrightarrow\mathfrak J_{B/A}\longrightarrow B\otimes_A B
 \xrightarrow p B\longrightarrow0
\]
分裂，与任意 $A'$ 张量后仍正合。又因为张量积与商相容，且理想平方的像等于像理想的平方，得到主部代数的基变换同构。最后用 $P^1_{B/A}=B\oplus\Omega^1_{B/A}$ 的分裂，取出微分模部分。正文已将这一理由写明。

\lead{N05：乘积代数的交叉分量。}
若 $B=B_1\times B_2$，令 $e_1,e_2$ 为对应正交幂等元。$e_1\otimes e_2$、$e_2\otimes e_1$ 在乘法映射下为零，且本身幂等，故属于 $\mathfrak J^2$。整个交叉分量由它们支撑，因而在一阶主部中消失。其余两分量正好给出 $P^1_{B_1/A}$、$P^1_{B_2/A}$；这是20.4.13(iii)的补明步骤。

\section*{底本、外部核验及检查范围}
\addcontentsline{toc}{section}{底本、外部核验及检查范围}
本轮翻译逐段依据上传法文原页；公式、上下标和交换图与页面图像对照，不把自动提取中丢失的上标当成作者错误。第二底本为IV第四部分的作者勘误。微分模与导子的泛性质另外参照Stacks Project的Tag 00RM、00RO作交叉核验；网页内容没有替代原文翻译，也没有被冒充作者勘误。

工程保存新增条目与公式标签清单、旧文件保留校验、有限例子验证脚本以及实际编译和渲染检查结果。文件保留校验仅说明哪些既有内容没有变动，不证明这些内容已无逻辑问题。单文件与分文件构建一致性，仅说明交付两种源码产生同一版面，不是独立数学审稿。

当前第0章已经连续处理至20.5.15；20.6--20.7、§§21--23、第IV章主体及第V卷预备笔记余文尚未译完。历史H8-P01等待核事项继续保留；本轮没有将这些问题宣布解决。
